% GENERATED FILE. Edit canonical lessons, then run npm run generate:books.
% canonical-source-hash: fnv1a64:f471f8236427f0b6
% canonical-lessons: GU-C338-granth, GU-C338-samayik, GU-C338-chapavvu, GU-C338-bhajavvu, GU-C338-adbhut, GU-R338-review-332-338, GU-R338-review-final-338

\chapter{A Book Comes Out}
\label{ch:gu-a-book-comes-out}
\hlchaptermodality{338}

\begin{chapteropening}
\textbf{By the end of this chapter:} \emph{I can say how a book was published and what people thought of it.}

Complete the last lesson of chapter 338: I can say how a book was published and what people thought of it.
\end{chapteropening}

\section[\texorpdfstring{\gu{ગ્રંથ} (granth)}{ગ્રંથ (granth)}]{\gu{ગ્રંથ} (granth) — a tome; a great book}

\label{lesson:GU-C338-granth}



\begin{quote}
Before the new one: say the Gujarati for the hero (of a story), then the Gujarati for a reader.
\end{quote}

\subsection*{You'll want to know: \gu{ગ્રંથ}}
\textbf{\gu{ગ્રંથ}} — \emph{granth} — \textquotedblleft{}a tome; a great book\textquotedblright{}.

\begin{quote}
\gu{પુસ્તકાલયમાં} \gu{એક} \gu{પ્રાચીન} \gu{ગ્રંથ} \gu{છે}.
\end{quote}

\begin{quote}
\emph{pustakālaymā̃ ek prācīn granth che.}
\end{quote}

\begin{quote}
There is an ancient tome in the library.
\end{quote}

A masculine noun.

\begin{cousinweb}
\textbf{\gu{ગ્રંથ}} (\emph{granth}) is a tatsama, Sanskrit \emph{grantha}, \textquotedblleft{}a knot\textquotedblright{}, and then \textquotedblleft{}a book\textquotedblright{}, once leaves tied together (CDIAL 4350). It is a large or important book.
\end{cousinweb}

\begin{grammarlens}[title={What you've built}]
One more word for telling what happened.
\end{grammarlens}

\subsection*{Your turn}
\begin{itemize}
\raggedright
  \item \emph{Say it:} \emph{granth}
  \item \emph{Say it:} \emph{granth}, once more
  \item \emph{Say it:} \emph{pustakālaymā̃ ek prācīn granth che.} — There is an ancient tome in the library.
  \item \emph{Recall:} say the Gujarati for the hero (of a story), then the Gujarati for a reader, then say \emph{granth} again
\end{itemize}

\begin{tcolorbox}[breakable,colback=teal!4,colframe=teal!35!black,title={Before you move on}]
What does \gu{ગ્રંથ} mean? (\textquotedblleft{}A tome; a great book\textquotedblright{}.) Say it once more.
\end{tcolorbox}
\section[\texorpdfstring{\gu{સામયિક} (sāmayik)}{સામયિક (sāmayik)}]{\gu{સામયિક} (sāmayik) — a magazine}

\label{lesson:GU-C338-samayik}



\begin{quote}
Before the new one: say the Gujarati for a reader, then the Gujarati for a tome.
\end{quote}

\subsection*{You'll want to know: \gu{સામયિક}}
\textbf{\gu{સામયિક}} — \emph{sāmayik} — \textquotedblleft{}a magazine\textquotedblright{}.

\begin{quote}
\gu{આ} \gu{સામયિક} \gu{દરેક} \gu{મહિને} \gu{આવે} \gu{છે}.
\end{quote}

\begin{quote}
\emph{ā sāmayik darek mahine āve che.}
\end{quote}

\begin{quote}
This magazine comes out every month.
\end{quote}

A neuter noun.

\begin{cousinweb}
\textbf{\gu{સામયિક}} (\emph{sāmayik}) is a tatsama, Sanskrit \emph{sāmayika}, \textquotedblleft{}seasonable, timely\textquotedblright{} (Monier-Williams); the sense of a magazine that comes out at set times is the modern language's.
\end{cousinweb}

\begin{grammarlens}[title={What you've built}]
One more word for telling what happened.
\end{grammarlens}

\subsection*{Your turn}
\begin{itemize}
\raggedright
  \item \emph{Say it:} \emph{sāmayik}
  \item \emph{Say it:} \emph{sāmayik}, once more
  \item \emph{Say it:} \emph{ā sāmayik darek mahine āve che.} — This magazine comes out every month.
  \item \emph{Recall:} say the Gujarati for a reader, then the Gujarati for a tome, then say \emph{sāmayik} again
\end{itemize}

\begin{tcolorbox}[breakable,colback=teal!4,colframe=teal!35!black,title={Before you move on}]
What does \gu{સામયિક} mean? (\textquotedblleft{}A magazine\textquotedblright{}.) Say it once more.
\end{tcolorbox}
\section[\texorpdfstring{\gu{છપાવવું} (chapāvvũ)}{છપાવવું (chapāvvũ)}]{\gu{છપાવવું} (chapāvvũ) — to get (something) printed; to publish}

\label{lesson:GU-C338-chapavvu}



\begin{quote}
Before the new one: say the Gujarati for a tome, then the Gujarati for a magazine.
\end{quote}

\subsection*{You'll want to know: \gu{છપાવવું}}
\textbf{\gu{છપાવવું}} — \emph{chapāvvũ} — \textquotedblleft{}to get (something) printed; to publish\textquotedblright{}.

\begin{quote}
\gu{લેખકે} \gu{પોતાની} \gu{નવલકથા} \gu{છપાવી}.
\end{quote}

\begin{quote}
\emph{lekhke potānī navalkathā chapāvī.}
\end{quote}

\begin{quote}
The author had his novel published.
\end{quote}

Past, with \gu{તેણે} (he or she …): \gu{તેણે} \gu{પુસ્તક} \gu{છપાવ્યું}, \emph{teṇe pustak chapāvyũ}.

\begin{cousinweb}
\textbf{\gu{છપાવવું}} (\emph{chapāvvũ}) is the causative of \textbf{\gu{છાપવું}}, to print, from \emph{*chapp}, \textquotedblleft{}press, cover\textquotedblright{}, a form CDIAL reconstructs (CDIAL 4994): to have a printer print it. The long \emph{ā} of \textbf{\gu{છાપવું}} shortens in the longer verb.
\end{cousinweb}

\begin{grammarlens}[title={What you've built}]
One more word for telling what happened.
\end{grammarlens}

\subsection*{Your turn}
\begin{itemize}
\raggedright
  \item \emph{Say it:} \emph{chapāvvũ}
  \item \emph{Say it:} \emph{chapāvvũ}, once more
  \item \emph{Say it:} \emph{lekhke potānī navalkathā chapāvī.} — The author had his novel published.
  \item \emph{Recall:} say the Gujarati for a tome, then the Gujarati for a magazine, then say \emph{chapāvvũ} again
\end{itemize}

\begin{tcolorbox}[breakable,colback=teal!4,colframe=teal!35!black,title={Before you move on}]
What does \gu{છપાવવું} mean? (\textquotedblleft{}To get (something) printed; to publish\textquotedblright{}.) Say it once more.
\end{tcolorbox}
\section[\texorpdfstring{\gu{ભજવવું} (bhajavvũ)}{ભજવવું (bhajavvũ)}]{\gu{ભજવવું} (bhajavvũ) — to perform (a play or a part)}

\label{lesson:GU-C338-bhajavvu}



\begin{quote}
Before the new one: say the Gujarati for a magazine, then the Gujarati for to publish.
\end{quote}

\subsection*{You'll want to know: \gu{ભજવવું}}
\textbf{\gu{ભજવવું}} — \emph{bhajavvũ} — \textquotedblleft{}to perform (a play or a part)\textquotedblright{}.

\begin{quote}
\gu{અરુણે} \gu{નાટકમાં} \gu{રાજાનું} \gu{પાત્ર} \gu{ભજવ્યું}.
\end{quote}

\begin{quote}
\emph{Aruṇe nāṭakmā̃ rājānũ pātra bhajavyũ.}
\end{quote}

\begin{quote}
Arun played the king in the play.
\end{quote}

Past, with \gu{તેણે} (he or she …): \gu{તેણે} \gu{નાટક} \gu{ભજવ્યું}, \emph{teṇe nāṭak bhajavyũ}.

\begin{cousinweb}
\textbf{\gu{ભજવવું}} (\emph{bhajavvũ}) No dictionary we use records an older form, so its history is not settled. Actors \textbf{\gu{નાટક} \gu{ભજવે} \gu{છે}}, put on a play, and each \textbf{\gu{પાત્ર} \gu{ભજવે} \gu{છે}}, plays a part.
\end{cousinweb}

\begin{grammarlens}[title={What you've built}]
One more word for telling what happened.
\end{grammarlens}

\subsection*{Your turn}
\begin{itemize}
\raggedright
  \item \emph{Say it:} \emph{bhajavvũ}
  \item \emph{Say it:} \emph{bhajavvũ}, once more
  \item \emph{Say it:} \emph{Aruṇe nāṭakmā̃ rājānũ pātra bhajavyũ.} — Arun played the king in the play.
  \item \emph{Recall:} say the Gujarati for a magazine, then the Gujarati for to publish, then say \emph{bhajavvũ} again
\end{itemize}

\begin{tcolorbox}[breakable,colback=teal!4,colframe=teal!35!black,title={Before you move on}]
What does \gu{ભજવવું} mean? (\textquotedblleft{}To perform (a play or a part)\textquotedblright{}.) Say it once more.
\end{tcolorbox}
\section[\texorpdfstring{\gu{અદ્ભુત} (adbhut)}{અદ્ભુત (adbhut)}]{\gu{અદ્ભુત} (adbhut) — wonderful}

\label{lesson:GU-C338-adbhut}



\begin{quote}
Before the new one: say the Gujarati for to publish, then the Gujarati for to perform.
\end{quote}

\subsection*{You'll want to know: \gu{અદ્ભુત}}
\textbf{\gu{અદ્ભુત}} — \emph{adbhut} — \textquotedblleft{}wonderful\textquotedblright{}.

\begin{quote}
\gu{આ} \gu{નવલકથા} \gu{અદ્ભુત} \gu{છે}.
\end{quote}

\begin{quote}
\emph{ā navalkathā adbhut che.}
\end{quote}

\begin{quote}
This novel is wonderful.
\end{quote}

\begin{cousinweb}
\textbf{\gu{અદ્ભુત}} (\emph{adbhut}) is a tatsama, Sanskrit \emph{adbhuta}, \textquotedblleft{}wonderful, marvellous\textquotedblright{} (CDIAL 240). It does not change for gender: \textbf{\gu{અદ્ભુત} \gu{નવલકથા}}, \textbf{\gu{અદ્ભુત} \gu{ગીત}}.
\end{cousinweb}

\begin{grammarlens}[title={What you've built}]
One more word for telling what happened.
\end{grammarlens}

\subsection*{Your turn}
\begin{itemize}
\raggedright
  \item \emph{Say it:} \emph{adbhut}
  \item \emph{Say it:} \emph{adbhut}, once more
  \item \emph{Say it:} \emph{ā navalkathā adbhut che.} — This novel is wonderful.
  \item \emph{Recall:} say the Gujarati for to publish, then the Gujarati for to perform, then say \emph{adbhut} again
\end{itemize}

\begin{tcolorbox}[breakable,colback=teal!4,colframe=teal!35!black,title={Before you move on}]
What does \gu{અદ્ભુત} mean? (\textquotedblleft{}Wonderful\textquotedblright{}.) Say it once more.
\end{tcolorbox}
\section[(dialogue)]{Looking back: from Once Upon a Time to A Book Comes Out}

\label{lesson:GU-R338-review-332-338}



\begin{quote}
Say the Gujarati for a legend and god.
\end{quote}

\subsection*{Your turn}
Say these aloud:

\begin{itemize}
\raggedright
  \item a legend, god, a demon, a goddess, a boon
  \item a fairy, magic, a miracle, vanished, to appear
  \item a treasure, a python, greed, to tempt, to wander
  \item a holy man, a curse, a moral, a proverb, to recite
  \item an author, a poet, to compose, an autobiography, imagination
  \item a page, a chapter (of a book), a character, the hero (of a story), a reader
  \item a tome, a magazine, to publish, to perform, wonderful
\end{itemize}

\begin{tcolorbox}[breakable,colback=teal!4,colframe=teal!35!black,title={Before you move on}]
A legend? (\textbf{\gu{દંતકથા}}, \emph{dantakthā}.) God? (\textbf{\gu{ભગવાન}}, \emph{bhagvān}.) A demon? (\textbf{\gu{દૈત્ય}}, \emph{daitya}.) A goddess? (\textbf{\gu{દેવી}}, \emph{devī}.) A boon? (\textbf{\gu{વરદાન}}, \emph{vardān}.) A fairy? (\textbf{\gu{પરી}}, \emph{parī}.) Magic? (\textbf{\gu{જાદુ}}, \emph{jādu}.) A miracle? (\textbf{\gu{ચમત્કાર}}, \emph{camatkār}.) Vanished? (\textbf{\gu{ગાયબ}}, \emph{gāyab}.) To appear? (\textbf{\gu{દેખાવું}}, \emph{dekhāvũ}.) A treasure? (\textbf{\gu{ખજાનો}}, \emph{khajāno}.) A python? (\textbf{\gu{અજગર}}, \emph{ajgar}.) Greed? (\textbf{\gu{લોભ}}, \emph{lobh}.) To tempt? (\textbf{\gu{લલચાવવું}}, \emph{lalcāvvũ}.) To wander? (\textbf{\gu{ભટકવું}}, \emph{bhaṭakvũ}.) A holy man? (\textbf{\gu{સાધુ}}, \emph{sādhu}.) A curse? (\textbf{\gu{શ્રાપ}}, \emph{śrāp}.) A moral? (\textbf{\gu{બોધ}}, \emph{bodh}.) A proverb? (\textbf{\gu{કહેવત}}, \emph{kahevat}.) To recite? (\textbf{\gu{સંભળાવવું}}, \emph{sambhaḷāvvũ}.) An author? (\textbf{\gu{લેખક}}, \emph{lekhak}.) A poet? (\textbf{\gu{કવિ}}, \emph{kavi}.) To compose? (\textbf{\gu{રચવું}}, \emph{racvũ}.) An autobiography? (\textbf{\gu{આત્મકથા}}, \emph{ātmakthā}.) Imagination? (\textbf{\gu{કલ્પના}}, \emph{kalpanā}.) A page? (\textbf{\gu{પાનું}}, \emph{pānũ}.) A chapter (of a book)? (\textbf{\gu{પ્રકરણ}}, \emph{prakaraṇ}.) A character? (\textbf{\gu{પાત્ર}}, \emph{pātra}.) The hero (of a story)? (\textbf{\gu{નાયક}}, \emph{nāyak}.) A reader? (\textbf{\gu{વાચક}}, \emph{vācak}.) A tome? (\textbf{\gu{ગ્રંથ}}, \emph{granth}.) A magazine? (\textbf{\gu{સામયિક}}, \emph{sāmayik}.) To publish? (\textbf{\gu{છપાવવું}}, \emph{chapāvvũ}.) To perform? (\textbf{\gu{ભજવવું}}, \emph{bhajavvũ}.) Wonderful? (\textbf{\gu{અદ્ભુત}}, \emph{adbhut}.)
\end{tcolorbox}
\section[(dialogue)]{A Book Comes Out — once more}

\label{lesson:GU-R338-review-final-338}



\begin{quote}
Say the Gujarati for a tome and a magazine.
\end{quote}

\subsection*{Your turn}
\begin{itemize}
\raggedright
  \item \emph{Say it:} a tome, a magazine, to publish, to perform, wonderful
\end{itemize}

\begin{tcolorbox}[breakable,colback=teal!4,colframe=teal!35!black,title={Before you move on}]
A tome? (\textbf{\gu{ગ્રંથ}}, \emph{granth}.) A magazine? (\textbf{\gu{સામયિક}}, \emph{sāmayik}.) To publish? (\textbf{\gu{છપાવવું}}, \emph{chapāvvũ}.) To perform? (\textbf{\gu{ભજવવું}}, \emph{bhajavvũ}.) Wonderful? (\textbf{\gu{અદ્ભુત}}, \emph{adbhut}.)
\end{tcolorbox}
